\section{Orientation: purpose, conventions and how to read this document}

\subsection{What this project is}
The project builds driving policies for the MetaDrive simulator that act from a \emph{bird's-eye view} (a top-down 64$\times$64 grid around the car, abbreviated BEV in the
code and in the teammates' documents, but written out here) plus a few scalar quantities, and asks how well such policies drive in traffic and when the
perception that feeds the grid is degraded by simulated fog and rain. The repository has three layers of work:
\begin{enumerate}
\item \textbf{Teammates' modules} (on the public GitHub repository): a vector-observation behaviour-cloning baseline, a five-camera rig with stereo depth, a YOLO object-detector with
  an image-space weather filter, the assembly of detections into the bird's-eye-view grid, and a behaviour-cloning policy on that grid.
\item \textbf{Policy-side work done in this project} (so far only on this machine, not on GitHub): fast camera-free environments, teacher policies, dataset-aggregation training,
  reinforcement-learning experiments, evaluation harnesses, perception measurements, diagnostics and recordings.
\item \textbf{This report} and the clean-up of the repository around it (Section~\ref{sec:cleanup}).
\end{enumerate}

\subsection{Naming conventions used in this document}
Names in the code were chosen quickly during experiments and are often abbreviations. In this document and in every figure and table, models and methods are
called by \emph{descriptive names}. Wherever a model also has a short code name (the name of its folder or of its result file), that code name is shown
once in typewriter font so that it can be found in the files. Table~\ref{tab:glossary} translates every abbreviation and code name that appears in the files.

\begin{table}[H]\centering\small
\caption{Glossary: abbreviations and code names used in file names, in the code and in earlier notes, with the full descriptive name used in this report.}
\label{tab:glossary}
\begin{tabularx}{\textwidth}{@{}>{\raggedright\arraybackslash}p{3.6cm}X@{}}\toprule
Code name or abbreviation & Meaning in this report\\\midrule
BEV & Bird's-eye view: a 64$\times$64 grid, 20\,m wide and 20\,m long (17.5\,m ahead of the car, 2.5\,m behind), cell size 0.3125\,m, centred on the ego car.\\
IDM / stock IDM & Intelligent driver model: MetaDrive's built-in rule-based car-following controller. Used as a teacher; it sees exact simulator positions of all vehicles.\\
SafeTeacher / \code{safe_teacher} & The collision-aware teacher: the intelligent driver model with lane changes switched off plus a predictive braking layer (Section~\ref{sec:teachers}).\\
BC / behaviour cloning & Supervised imitation of recorded expert actions (regression of steering and throttle).\\
DAgger / dataset aggregation & Iterative imitation: the learner drives, the teacher labels the states the learner visited, and the labelled states are added to the training set (Ross et al., 2011).\\
PPO / proximal policy optimisation & The reinforcement-learning algorithm used in every reinforcement-learning attempt.\\
HD map & High-definition map: here the lane centre line and the road boundaries read from the simulator's map; treated as a legitimate structured input.\\
\code{repo_bc_clear} & Baseline behaviour cloning, trained on five demonstrations \emph{without} traffic (\pth{bc_clear_5eps_without_traffic.pt}).\\
\code{repo_bc_sparse} & Baseline behaviour cloning, trained on five demonstrations \emph{with} sparse traffic (\pth{bc_sparse_traffic_5eps.pt}). The repository calls it the recommended starting checkpoint.\\
\code{base_a} & Dataset-aggregation student on the teammates' unmodified one-channel view and all six scalars (no map overlay).\\
\code{ov_both}, \code{ov_bev_only} & Map-overlay view with all six scalars / with no scalars. (\code{ov} = overlay.)\\
\code{ov_split} & \textbf{The headline model}: map-overlay view and three scalars (speed, route progress, distance to goal); the three lane-geometry quantities were moved from the scalars into the overlay.\\
\code{arch_bev_only}, \code{arch_scalar_only} & Input-ablation models: view only / scalars only (\code{arch} = architecture ablation).\\
\code{caut_a}, \code{ctrl_nocaut}, \code{caut_c0}, \code{caut_auto} & Caution study: caution-conditioned network; its control with the caution input always zero; the same conditioned network evaluated with caution 0 and with caution derived from weather.\\
\code{deg_base} & Dataset-aggregation student trained under the calibrated degraded perception.\\
\code{ppo_plain}, \code{ppo_ada} & Reinforcement learning from \code{deg_base} without / with a weather input (\code{ada} = adaptive; this run did not finish).\\
\code{rl_plain}, \code{rl_energy}, \code{rl_safe}, \code{rl_head}, \code{rl_full} & Reward-combination sweep, reinforcement learning from the headline model (base reward; plus energy; plus safe-speed and scaled crash penalty; plus headway; all terms).\\
\code{C2}, \code{C3}, \code{probe_C}, \code{pilot_cond} & Early reinforcement-learning attempts (privileged teacher trained with PPO).\\
\code{pose_aug_*}, \code{pose_cam_*} & Pose-error study: map-pose-error augmentation (and two-cue variant with a noisy lane detector), conditioned or not on the error level.\\
\code{vcam} & Pixel policy: a small convolutional network that sees only a front-camera image.\\
Seed / seed block & Integer that selects the MetaDrive map, route and initial traffic. Disjoint blocks are used for training, validation and testing (Table~\ref{tab:seeds}).\\
Density & MetaDrive \code{traffic_density}: probability-like parameter controlling the amount of traffic. 0.05 = sparse, 0.2 = dense.\\
Wh/km & Estimated traction energy per kilometre from the declared virtual-vehicle model in \pth{bev_rl/energy.py}; not battery electricity.\\\bottomrule
\end{tabularx}
\end{table}

\subsection{Evidence levels}
Every quantitative statement in the document carries its sample size. Three levels are used:
\textbf{reported} ($n\ge 200$ paired episodes, confidence intervals and paired tests given),
\textbf{pilot} ($n=10$ to $40$ episodes; wide intervals; used to choose what to run next, not to claim),
\textbf{diagnostic} (one or two replayed episodes that explain a mechanism but do not measure a rate).

\subsection{How to read the document}
Sections~\ref{sec:repomap}--\ref{sec:perception} explain the setting; Sections~\ref{sec:arch}--\ref{sec:catalogue} describe the models; Sections~\ref{sec:headline}--\ref{sec:failure}
give the results and failures; Section~\ref{sec:videos} describes the videos; Section~\ref{sec:lessons} summarises what worked and what did not; Section~\ref{sec:directions} lists
paper directions and the evidence behind each; Section~\ref{sec:related} lists related 2025--2026 papers. Appendix~\ref{app:inventory} lists every file created or changed, and
Appendix~\ref{app:honesty} lists known defects and invalid runs.
